% !TEX program = xelatex
% NOTE: This manuscript is typeset for XeLaTeX.

\documentclass[12pt,a4paper]{article}

\usepackage{fontspec}
\setmainfont{TeX Gyre Termes}
\usepackage{xeCJK}
\setCJKmainfont{Noto Sans CJK SC}[AutoFakeBold]

\usepackage[a4paper,margin=2.5cm]{geometry}
\usepackage{microtype}
\usepackage[hidelinks]{hyperref}
\usepackage{url}
\usepackage{amsmath,amssymb}
\usepackage{graphicx}
\usepackage{booktabs}
\usepackage{tabularx}
\usepackage{multirow}
\usepackage{float}
\usepackage{enumitem}
\usepackage{xcolor}
\usepackage{caption}
\usepackage{setspace}
\usepackage{titlesec}
\usepackage{fancyhdr}
\usepackage{algorithm}
\usepackage{algpseudocode}

\setlength{\parindent}{1.5em}
\setlength{\parskip}{0.35em}
\setstretch{1.15}

\captionsetup{font=small,labelfont=bf,labelsep=quad,justification=centering}
\titleformat{\section}{\Large\bfseries}{\thesection}{0.75em}{}
\titleformat{\subsection}{\large\bfseries}{\thesubsection}{0.75em}{}
\titleformat{\subsubsection}{\normalsize\bfseries}{\thesubsubsection}{0.75em}{}

\pagestyle{fancy}
\fancyhf{}
\fancyhead[C]{\small Manuscript}
\fancyfoot[C]{\small \thepage}
\renewcommand{\headrulewidth}{0.4pt}

\newcolumntype{Y}{>{\centering\arraybackslash}X}
\newcolumntype{L}{>{\raggedright\arraybackslash}X}
\renewcommand{\arraystretch}{1.25}

\newcommand{\algcomment}[1]{\hfill\textcolor{gray}{\small // #1}}

% Simple figure placeholder
\newcommand{\figplaceholder}[2]{%
  \fbox{\parbox[c][0.28\textheight][c]{0.92\linewidth}{\centering\textbf{Figure Placeholder}\\#1\\\small (#2)}}%
}

\title{SIGNPOST-Bench: Evaluating Textual Adversarial Robustness for Visual Geo-localization with Multimodal Large Models}
\author{Sirun Li\\School of Earth and Space Sciences}
\date{March 2026}

\begin{document}
\maketitle

% ===========================================================================
%  ABSTRACT
% ===========================================================================
\begin{abstract}
Multimodal large language models (MLLMs) have demonstrated remarkable zero-shot visual geo-localization capabilities by integrating scene understanding with world knowledge. However, these models frequently exhibit a critical vulnerability: excessive reliance on embedded scene text (OCR bias), where textual cues override global visual evidence such as vegetation, geomorphology, and architectural style. To systematically quantify this failure mode, we propose \textbf{SIGNPOST-Bench} (\textbf{S}cene \textbf{I}mage \textbf{G}eo-localization with \textbf{N}oisy \textbf{P}erturbation on \textbf{O}bserved \textbf{S}ign \textbf{T}ext), a fully automated adversarial evaluation framework.

SIGNPOST-Bench introduces an \emph{LLM-in-the-loop} pipeline that operates in four stages: (1)~filtering images containing legible text via OCR pre-screening, (2)~employing a frontier VLM to analyze native scene text and automatically generate three categories of semantically controlled replacement text---\emph{Similar}, \emph{Random}, and \emph{Adversarial}, (3)~photorealistically injecting the replacement text via ComfyUI-orchestrated image editing, and (4)~evaluating target models on both clean baselines (text removed) and adversarial variants. This counterfactual design isolates the causal effect of textual interference from confounding visual factors.

We evaluate four representative MLLMs spanning 7B--30B parameters and two reasoning paradigms (chain-of-thought vs.\ instruction-following) on three geographically diverse datasets: YFCC4K, Im2GPS3K, and Google Street View. Results are reported using three complementary metrics: Weighted Localization Accuracy (WLA), Text Bias Score (TBS), and Trap-Fit Rate (TFR). The code and data are available at \url{https://github.com/inorganicwriter/SIGNPOST-Bench}.
\end{abstract}

\noindent\textbf{Keywords:} visual geo-localization; multimodal large language models; adversarial robustness; typographic attack; OCR bias; benchmark

% ===========================================================================
%  1. INTRODUCTION
% ===========================================================================
\section{Introduction}

Visual geo-localization---inferring the geographic coordinates of an image from its visual content---has evolved from a retrieval-based matching problem~\cite{im2gps,geoclip,pigeon} into a generative reasoning task driven by multimodal large language models (MLLMs). Modern MLLMs such as GPT-4V~\cite{gpt4v}, Qwen-VL~\cite{qwenvl}, and Gemini can leverage scene text (road signs, storefront names, address plates) alongside internal world knowledge to achieve impressive zero-shot localization accuracy.

However, this capability introduces a critical failure mode: \emph{modality dominance imbalance}. Because embedded text carries explicit geographic pointers (e.g., a street name or country label), models tend to develop a heuristic shortcut that prioritizes optical character recognition (OCR) output over holistic visual evidence---a phenomenon we term \emph{OCR bias}. When global context (vegetation patterns, architectural style, road infrastructure) conflicts with embedded text, models frequently ignore the former and over-trust the latter, leading to severe localization errors.

\begin{figure}[H]
  \centering
  \figplaceholder{OCR bias mechanism illustration}{Left: balanced attention without prominent text; Right: attention collapse to text regions under interference}
  \caption{Schematic of OCR bias in visual geo-localization. (a) Without prominent text, the model distributes attention across architectural style, vegetation, and road markings. (b) Under adversarial text interference, attention collapses to the text region, overriding visual evidence.}
  \label{fig:ocr_bias_mechanism}
\end{figure}

This vulnerability relates to typographic attacks on image classifiers~\cite{typographic_attack,azuma2023defense,qraitem2024typo}. However, while prior studies focus on class-level confusion (e.g., misclassifying an ``apple'' as an ``iPod''), SIGNPOST-Bench targets a spatial reasoning task where the attack \emph{redirects} the prediction to a specific geographical coordinate. Despite growing interest in MLLM robustness~\cite{hallusionbench,pope,naturalbench,mmrobustness}, the community lacks a benchmark that systematically evaluates text-induced bias in geo-localization with a rigorous counterfactual design.

We introduce SIGNPOST-Bench with the following contributions:

\begin{enumerate}[leftmargin=2em]
  \item \textbf{A fully automated, LLM-driven adversarial pipeline.} We leverage a frontier VLM (Qwen3-VL-30B) to analyze scene text and generate three categories of replacement text---\emph{Similar}, \emph{Random}, and \emph{Adversarial}---spanning a semantic gradient from near-identical perturbation to deliberate geographic misdirection. This eliminates the need for manual distractor curation and produces context-aware, image-specific attacks.

  \item \textbf{A two-dimensional evaluation framework.} We introduce a \emph{Scene-Text Coupling Taxonomy} that classifies native text by its geographic informativeness into three tiers---Portable (T1), Culturally Indicative (T2), and Geo-Specific (T3). Crossing these with the three attack categories yields a $3 \times 3$ evaluation matrix, enabling fine-grained analysis of how both the inherent informativeness of scene text and the semantic distance of the replacement jointly determine model vulnerability.

  \item \textbf{A counterfactual evaluation design.} Using ComfyUI~\cite{comfyui} with the Qwen-Image-Edit model, we construct paired image sets from the same source: a Blank baseline (text removed) and three adversarial variants (text replaced). By comparing model predictions on Blank vs.\ adversarial images, we obtain causal estimates of text-induced bias, free from confounds.

  \item \textbf{Multi-dimensional robustness metrics.} We define Weighted Localization Accuracy (WLA) for multi-scale precision, Text Bias Score (TBS) for causal effect measurement, and Trap-Fit Rate (TFR) for directional attack success quantification, providing a finer-grained evaluation than aggregate error distance alone.
\end{enumerate}

\begin{figure}[H]
  \centering
  \figplaceholder{SIGNPOST-Bench pipeline overview}{4-stage flow: OCR Filter → LLM Attack Gen → ComfyUI Synthesis → Multi-Model Evaluation}
  \caption{Overall architecture of SIGNPOST-Bench. The four-stage pipeline: (1) OCR-based text filtering, (2) LLM-driven attack generation, (3) ComfyUI image synthesis, and (4) multi-model evaluation with robustness metrics.}
  \label{fig:pipeline_overview}
\end{figure}


% ===========================================================================
%  2. RELATED WORK
% ===========================================================================
\section{Related Work}

\subsection{Visual geo-localization}
Visual geo-localization has progressed from image retrieval~\cite{im2gps} to deep representation learning. Im2GPS~\cite{im2gps} pioneered the field by matching global image descriptors against geo-tagged databases. Subsequent work improved feature representations through region-level matching~\cite{pigeon} and contrastive alignment between images and geographic coordinates~\cite{geoclip}. StreetCLIP~\cite{streetclip} achieved zero-shot generalization by adapting CLIP to the geo-localization domain using diverse street-view data, while GeoReasoner~\cite{georeasoner} introduced chain-of-thought prompting for expert-level geographic reasoning with MLLMs. GEOBench-VLM~\cite{geobenchvlm} recently benchmarked VLMs on a range of geospatial tasks, finding that even the best-performing models achieve only moderate accuracy on geospatial understanding, underscoring the challenge of this domain.

A parallel line of work studies \emph{cross-view} geo-localization, where satellite and ground-level image pairs are matched. Recent robustness benchmarks for cross-view tasks (e.g., CVUSA-C, CVACT-C~\cite{crossview_robustness}) measure model resilience under natural corruptions such as weather variations and sensor noise. SIGNPOST-Bench differs from these efforts by focusing on the \emph{semantic-level} textual interference intrinsic to street-view scenes rather than pixel-level corruption.

\subsection{Typographic attacks on vision-language models}
Typographic attacks exploit the text-reading capabilities of VLMs by embedding misleading text into images. Goh et al.~\cite{typographic_attack} first demonstrated that CLIP's multimodal neurons respond to handwritten text overlaid on images, enabling class-level confusion. Subsequent work has systematized this vulnerability: Azuma et al.~\cite{azuma2023defense} proposed defense strategies through prompt engineering, while Qraitem et al.~\cite{qraitem2024typo} introduced the TypoD benchmark for evaluating VLM distractibility under typographic interference. More recent work has explored \emph{self-generated} attacks, where the VLM itself proposes adversarial text~\cite{selfgenerated_typo}, and \emph{scene-coherent} attacks that blend naturalistically into the visual context~\cite{scenecoherent_typo}. The SCAM benchmark~\cite{scam} provides a large-scale dataset of real-world typographic attack images.

SIGNPOST-Bench extends the typographic attack paradigm in two key ways: (1)~we replace \emph{in-situ} scene text rather than overlaying new text, producing more photorealistic adversarial samples; and (2)~we measure \emph{location-specific} bias (geographic displacement) rather than class-level confusion, with metrics that quantify both the magnitude and directionality of the attack's effect.

\subsection{MLLM robustness and hallucination evaluation}
A growing body of work evaluates MLLMs for robustness and reliability. HallusionBench~\cite{hallusionbench} tests reasoning under visual illusions and context manipulation. POPE~\cite{pope} measures object hallucination frequency. NaturalBench~\cite{naturalbench} evaluates models on naturally adversarial samples. MME~\cite{mme} and MMBench~\cite{mmbench} provide comprehensive multi-task evaluations, frequently including document or scene text understanding modules. These benchmarks collectively reveal that MLLMs are susceptible to various failure modes including over-reliance on language priors, visual grounding failures, and context misinterpretation.

SIGNPOST-Bench contributes to this ecosystem by providing a \emph{domain-specific} (geo-localization), \emph{intervention-based} (counterfactual pairing) benchmark with \emph{continuous} (rather than binary) outcome metrics, enabling fine-grained analysis of text--visual conflict in a safety-critical application domain.


% ===========================================================================
%  3. PROBLEM FORMULATION
% ===========================================================================
\section{Problem Formulation}

\subsection{Task definition}
Let $f: \mathcal{I} \to \mathbb{R}^2$ denote a multimodal geo-localization model that maps an image $I$ to predicted GPS coordinates $(\hat{\phi}, \hat{\lambda})$. Given an image $I_{raw}$ with ground-truth coordinates $(\phi_{gt}, \lambda_{gt})$, we define the localization error as the geodesic (Haversine) distance:
\begin{equation}
  \mathcal{D}\bigl(f(I),\, (\phi_{gt}, \lambda_{gt})\bigr) = 2R \cdot \arcsin\!\sqrt{\sin^2\!\frac{\Delta\phi}{2} + \cos\phi_{gt}\cos\hat\phi\,\sin^2\!\frac{\Delta\lambda}{2}}
  \label{eq:haversine}
\end{equation}
where $R = 6371$\,km is the Earth's radius.

\subsection{Threat model}
We consider a \emph{semantic text replacement} attack. The attacker has access to the input image but not to the target model's parameters. The attacker's goal is to modify the scene text in $I_{raw}$ such that the model's prediction $f(I_{adv})$ is displaced toward a target ``trap'' location $(\phi_T, \lambda_T)$ implied by the injected text.

Formally, given a text replacement function $\mathcal{T}$ that substitutes native text with adversarial text $t_{adv}$, the attack is:
\begin{equation}
  I_{adv} = \mathcal{T}(I_{raw}, t_{adv}), \quad \text{s.t.} \quad \mathcal{D}\bigl(f(I_{adv}),\, (\phi_T, \lambda_T)\bigr) < \epsilon
\end{equation}
where $\epsilon = 50$\,km defines the trap region. To isolate the textual intervention, we also construct:
\begin{equation}
  I_{blank} = \mathcal{T}(I_{raw}, \varnothing) \qquad \text{(text removal)}
\end{equation}

The \emph{Text Bias Score} then quantifies the causal effect:
\begin{equation}
  \text{TBS} = \mathcal{D}\bigl(f(I_{adv}),\, y_{gt}\bigr) - \mathcal{D}\bigl(f(I_{blank}),\, y_{gt}\bigr)
  \label{eq:tbs}
\end{equation}

\subsection{Attack categories}
We define three attack categories that span a gradient of semantic distance from the original text:
\begin{itemize}[leftmargin=2em]
  \item \textbf{Similar} ($t_{sim}$): a visually or semantically close variant of the original text (e.g., ``McDonald's'' $\to$ ``McDonalds''; a phone number changed by one digit). Tests sensitivity to subtle perturbations.
  \item \textbf{Random} ($t_{rnd}$): a completely unrelated word or phrase (e.g., a Chinese storefront name $\to$ ``Galaxy Mart''). Tests whether the model blindly follows any OCR output regardless of visual incongruence.
  \item \textbf{Adversarial} ($t_{adv}$): text that conveys a deliberately misleading geographic cue (e.g., a Tokyo address $\to$ ``Paris''; ``Stop'' $\to$ ``Go''). Maximizes semantic conflict and tests cross-modal verification ability.
\end{itemize}

\subsection{Scene-text coupling taxonomy}
Orthogonal to the attack categories (which characterize the \emph{replacement} text), we classify the \emph{original} scene text by its inherent geographic informativeness. The ordering principle is \textbf{text--location coupling strength}: how much localization information does the text intrinsically carry?

\begin{table}[H]
  \centering
  \caption{Scene-text coupling taxonomy. Each tier is defined by the geographic informativeness of the native text, with annotation criteria and expected robustness behavior.}
  \label{tab:coupling_taxonomy}
  \resizebox{\textwidth}{!}{
    \begin{tabular}{l p{3.5cm} p{4cm} p{2.5cm} p{4cm}}
      \toprule
      \textbf{Tier} & \textbf{Definition} & \textbf{Examples} & \textbf{Geo-signal} & \textbf{Expected behavior} \\
      \midrule
      \textbf{T1: Portable} & Text not inherently tied to any location; could appear on any continent & Global brand logos (FedEx, Coca-Cola), vehicle text, clothing labels, product names & None & Robust model should ignore; attacks should have minimal impact \\
      \midrule
      \textbf{T2: Cultural} & Text that narrows location to a broad region via script or cultural convention & Non-Latin script (中文, العربية), generic local business names, language-specific menus & Region / country & Visual--linguistic conflict testable; model should weigh architectural style over text \\
      \midrule
      \textbf{T3: Geo-Specific} & Text that directly identifies a specific geographic entity & Street names, city names, postal codes, named landmarks, specific addresses & City / street & Highest vulnerability; attack provides a plausible redirect target \\
      \bottomrule
    \end{tabular}
  }
\end{table}

\subsubsection{T1: Portable text (weak coupling)}
This tier encompasses text on \emph{mobile} or \emph{globally ubiquitous} carriers: vehicle liveries, clothing, international brand signage, and product labels. A FedEx truck or a Coca-Cola billboard could appear in virtually any country. Because such text carries no geographic signal, a robust model should not change its prediction when the text is modified.

\textbf{Annotation criterion}: Could this exact text plausibly appear unchanged on a different continent?  If yes $\to$ T1.

\textbf{Robustness hypothesis}: Attacks on T1 text should produce near-zero TBS. Models with significant TBS on T1 text exhibit \emph{undiscriminating OCR dependence}---they process any detected text as a location cue, even when the text is geographically uninformative.

\subsubsection{T2: Culturally indicative text (medium coupling)}
This tier covers text that narrows the location to a \emph{cultural or linguistic region} but does not pinpoint a specific place. Examples include writing-system-specific text (Chinese characters on a storefront, Arabic on a street sign), generic local business types (e.g., ``药店'' / ``pharmacie'' / ``Apotheke''), and culturally coded signage conventions. The text implies a rough geography (East Asia, Middle East, Francophone region) but not a precise location.

\textbf{Annotation criterion}: Does the text narrow the location to a specific language or cultural region, without identifying a specific city or address?  If yes $\to$ T2.

\textbf{Robustness hypothesis}: Replacing T2 text with text from a different cultural region (e.g., Chinese $\to$ English on a distinctly Chinese-style building) creates a \emph{visual--linguistic conflict}. A robust model should detect this inconsistency---the architectural style, road markings, and vegetation all contradict the injected text---and rely on the globally coherent visual context.

\subsubsection{T3: Geo-specific text (strong coupling)}
This tier addresses text that \emph{directly identifies} a geographic entity: street names, city labels, postal codes, named landmarks, and explicit addresses. Such text provides a strong, high-confidence location cue. Replacing a street name with an address from a different city creates a \emph{precise misdirection} scenario.

\textbf{Annotation criterion}: Can a human determine a specific city or sub-city location from the text alone?  If yes $\to$ T3.

\textbf{Robustness hypothesis}: T3 text presents the highest vulnerability. Because the text itself constitutes strong evidence, even models with good visual grounding may be swayed by an adversarial replacement---especially when the injected text is semantically plausible (Adversarial category). The combination of T3 $\times$ Adversarial represents the most challenging scenario in our evaluation matrix.

\begin{figure}[H]
  \centering
  \figplaceholder{3x3 Evaluation Matrix Diagram}{Visual grid showing Tier (T1, T2, T3) vs Attack Category (Similar, Random, Adversarial) with example perturbations.}
  \caption{The two-dimensional evaluation matrix of SIGNPOST-Bench, crossing original text coupling strength with attacked text semantic deviation.}
  \label{fig:3x3_matrix}
\end{figure}

\vspace{0.5em}
\noindent The coupling taxonomy is applied via \textbf{manual annotation} on a representative subset of benchmark images. The annotator inspects the native scene text identified by the VLM in Stage~2 and assigns a tier label based on the criteria above. This annotation enables a $3 \times 3$ analysis matrix (Tier $\times$ Attack Category) that reveals how the geographic informativeness of the original text modulates the effectiveness of different attack strategies.


% ===========================================================================
%  4. THE SIGNPOST-BENCH FRAMEWORK
% ===========================================================================
\section{The SIGNPOST-Bench Framework}
SIGNPOST-Bench implements a fully automated, four-stage pipeline (Figure~\ref{fig:pipeline_overview}). Two external services underpin the system: a frontier VLM (Qwen3-VL-30B-A3B-Thinking~\cite{qwen3vl}) deployed via vLLM~\cite{vllm} for text analysis and model evaluation, and ComfyUI~\cite{comfyui} with the Qwen-Image-Edit model for photorealistic image editing.

\subsection{Stage 1: Data collection and text filtering}
To ensure the benchmark contains meaningful textual content, we first filter raw images using optical character recognition. We employ EasyOCR to detect legible text in each candidate image. Images with fewer than two detected characters are discarded, ensuring that only images where text is genuinely present---and potentially influential for geo-localization---enter the pipeline.

We source images from three datasets to ensure geographic diversity:
\begin{itemize}[leftmargin=2em]
  \item \textbf{YFCC4K}: a curated 4,000-image subset of YFCC100M~\cite{yfcc100m}, with ground-truth coordinates.
  \item \textbf{Im2GPS3K}: the established 3,000-image test set from the Im2GPS benchmark~\cite{im2gps}.
  \item \textbf{Google Street View}: a stratified sample drawn proportionally by country from a global panoramic dataset, covering underrepresented regions in the above datasets.
\end{itemize}

Ground-truth metadata (latitude, longitude) is converted to a unified format compatible with the evaluation pipeline.

\subsection{Stage 2: LLM-driven attack generation}
For each filtered image, we prompt the VLM with the raw image and instruct it to perform three tasks in a single inference call:

\begin{enumerate}[leftmargin=2em]
  \item \textbf{Identify} the main text content (e.g., store names, road signs, address plates).
  \item \textbf{Localize} where the text appears, described in natural language (e.g., ``on the green storefront sign at the top center'').
  \item \textbf{Generate} three replacement texts: Similar, Random, and Adversarial.
\end{enumerate}

The VLM outputs a structured JSON response, which is automatically parsed and validated. Entries where no legible text is detected---or where the JSON parsing fails---are filtered out. The attacks are saved to a JSONL file (\texttt{attacks.jsonl}), with each entry recording the original text, its spatial location description, and the three generated replacement texts.

\begin{algorithm}[H]
  \caption{LLM-Driven Attack Generation (\texttt{generate\_attacks.py})}
  \begin{algorithmic}[1]
    \Require Image set $\{I_1, \ldots, I_N\}$; VLM $\mathcal{M}$
    \Ensure Attack configurations $\mathcal{A} = \{a_1, \ldots, a_M\}$ \quad ($M \le N$)
    \For{each image $I_i$ \textbf{in parallel} (semaphore = 20)}
      \State $\mathit{response} \leftarrow \mathcal{M}(I_i, \mathit{prompt}_{attack})$ \algcomment{single VLM call per image}
      \State Parse $\mathit{response}$ as JSON $\to (text_i, loc_i, \{t_{sim}, t_{rnd}, t_{adv}\})$
      \If{parsing succeeds \textbf{and} $\text{attacks} \neq \varnothing$}
        \State Append $a_i = (I_i, text_i, loc_i, t_{sim}, t_{rnd}, t_{adv})$ to $\mathcal{A}$
      \EndIf
    \EndFor
    \State \Return $\mathcal{A}$
  \end{algorithmic}
\end{algorithm}

This \emph{LLM-in-the-loop} approach offers key advantages over prior static distractor pools~\cite{typographic_attack} or CLIP-based selection methods:
\begin{itemize}
  \item \textbf{Context-awareness}: attacks are tailored to each image's specific text content and scene context.
  \item \textbf{Scalability}: the pipeline processes arbitrary image collections without manual curation.
  \item \textbf{Semantic precision}: the three categories provide a principled gradient from minimal to maximal semantic conflict.
\end{itemize}

\subsection{Stage 3: Counterfactual image synthesis}
For each attack configuration, we synthesize images using ComfyUI~\cite{comfyui} with the Qwen-Image-Edit model. The synthesis produces four variants from each original image $I_{raw}$:

\begin{enumerate}[leftmargin=2em]
  \item $I_{blank}$\,: text removed \hfill \texttt{"Remove the text \{loc\}. Maintain photorealism."}
  \item $I_{sim}$\,: text replaced with Similar variant \hfill \texttt{"Replace the text \{loc\} with '\{$t_{sim}$\}'."}
  \item $I_{rnd}$\,: text replaced with Random variant
  \item $I_{adv}$\,: text replaced with Adversarial variant
\end{enumerate}

The image editing workflow operates via WebSocket communication: images are uploaded to ComfyUI, processed through the Qwen-Image-Edit diffusion model with the constructed prompt, and the generated results are downloaded and saved with structured filenames encoding the attack type and injected text.

\begin{algorithm}[H]
  \caption{Counterfactual Image Synthesis (\texttt{main\_benchmark.py})}
  \begin{algorithmic}[1]
    \Require Attack set $\mathcal{A}$; Image editor $\mathcal{E}$ (ComfyUI + Qwen-Image-Edit)
    \Ensure Image quintuplets $\{(I_{raw}, I_{blank}, I_{sim}, I_{rnd}, I_{adv})\}$
    \For{each $(I_i, text_i, loc_i, t_{sim}, t_{rnd}, t_{adv})$ \textbf{in} $\mathcal{A}$}
      \State $I_{blank} \leftarrow \mathcal{E}(I_i,\, \texttt{"Remove the text } loc_i \texttt{."})$ \algcomment{control group}
      \For{$(t, label)$ in $\{(t_{sim}, sim), (t_{rnd}, rnd), (t_{adv}, adv)\}$}
        \State $I_{label} \leftarrow \mathcal{E}(I_i,\, \texttt{"Replace the text } loc_i \texttt{ with '} t \texttt{'."})$
      \EndFor
      \State Save all variants; append provenance to \texttt{benchmark\_meta.jsonl}
    \EndFor
  \end{algorithmic}
\end{algorithm}

All four variants share the same visual substrate, differing only in their textual content. The pipeline supports batch processing with automatic resume capability: existing output files are detected and skipped, enabling recovery from interruptions.

\begin{figure}[H]
  \centering
  \figplaceholder{Synthesis example: $I_{raw} \to I_{blank} \to I_{sim} \to I_{rnd} \to I_{adv}$}{showing photorealistic text replacement across all 4 variants}
  \caption{Image synthesis example. Left to right: raw image $I_{raw}$, Blank baseline $I_{blank}$ (text removed), Similar $I_{sim}$, Random $I_{rnd}$, and Adversarial $I_{adv}$ variants.}
  \label{fig:synthesis_example}
\end{figure}

\subsection{Stage 4: Multi-model evaluation}
The evaluation stage feeds synthesized images to target models and computes metrics. For each image, the model is prompted to output GPS coordinates, which are parsed using a cascading strategy: JSON extraction, regex matching of $(lat, lon)$ format, and labeled format (``Latitude: ..., Longitude: ...'').

The evaluation pipeline first processes original (unmodified) images to establish a baseline error for each sample, then evaluates each attack subdirectory (Similar, Random, Adversarial) separately. TBS is computed by pairing each adversarial result with its corresponding original-image baseline via filename matching.


% ===========================================================================
%  5. EVALUATION METRICS
% ===========================================================================
\section{Evaluation Metrics}
We define three complementary metrics (Table~\ref{tab:metrics_summary}).

\begin{table}[H]
  \centering
  \caption{Mathematical definitions of SIGNPOST-Bench evaluation metrics}
  \label{tab:metrics_summary}
  \resizebox{\textwidth}{!}{
    \begin{tabularx}{\textwidth}{l Y X}
      \toprule
      \textbf{Metric} & \textbf{Symbol} & \textbf{Definition} \\
      \midrule
      Weighted Localization Accuracy & WLA & $S_i = \sum_{k=1}^{5} w_k \cdot \mathbb{I}(d_i < \tau_k)$ \\
      \midrule
      Text Bias Score & TBS & $\text{TBS}_i = \mathcal{D}(\hat{y}_{adv}, y_{gt}) - \mathcal{D}(\hat{y}_{blank}, y_{gt})$ \\
      \midrule
      Trap-Fit Rate & TFR & $\text{TFR} = \frac{1}{N}\sum_{i=1}^{N} \mathbb{I}\!\left(\mathcal{D}(\hat{y}_{adv}^{(i)}, y_{trap}^{(i)}) < 50\,\text{km}\right)$ \\
      \bottomrule
    \end{tabularx}
  }
\end{table}

\noindent\textbf{Weighted Localization Accuracy (WLA).}
Geo-localization errors span orders of magnitude---from street-level (sub-kilometer) to continental (thousands of kilometers). To capture multi-scale performance, we define five distance thresholds $\tau_k \in \{1, 25, 200, 750, 2500\}$\,km, covering street, city, region, country, and continental scales. Each threshold carries equal weight $w_k = 0.2$:
\begin{equation}
  \text{WLA}_i = \sum_{k=1}^{5} 0.2 \cdot \mathbb{I}(d_i < \tau_k), \qquad d_i = \mathcal{D}(f(I_i), y_{gt}^{(i)})
\end{equation}
WLA $\in [0\%, 100\%]$; higher values indicate more precise localization. An average WLA of 100\% means every prediction is within 1\,km of ground truth. In tables, WLA is reported as a percentage (e.g., 57.5\%).

\noindent\textbf{Text Bias Score (TBS).}
TBS measures the \emph{causal effect} of textual interference by comparing the prediction error on an adversarial image against its paired Blank baseline (Eq.~\ref{eq:tbs}). TBS $> 0$ indicates that the injected text increased localization error; larger positive values indicate stronger text-induced bias. TBS can be computed per attack category to reveal differential vulnerability.

\noindent\textbf{Trap-Fit Rate (TFR).}
TFR quantifies the \emph{directional} success of adversarial attacks. It measures the proportion of predictions that fall within 50\,km of the location implied by the injected text (the ``trap region''). High TFR indicates that the model was not merely confused but was \emph{specifically redirected} to the attacker's intended location---a stronger indication of OCR bias than TBS alone.


% ===========================================================================
%  6. EXPERIMENTS AND RESULTS
% ===========================================================================
\section{Experiments and Results}

\subsection{Experimental setup}

\textbf{Models.} We evaluate four representative MLLMs spanning different parameter scales and reasoning paradigms (Table~\ref{tab:model_list}).

\begin{table}[H]
  \centering
  \caption{Models evaluated in SIGNPOST-Bench}
  \label{tab:model_list}
  \begin{tabular}{llll}
    \toprule
    \textbf{Model} & \textbf{Params} & \textbf{Reasoning Style} & \textbf{Deployment} \\
    \midrule
    Qwen3-VL-30B-A3B-Thinking & 30B & Chain-of-thought & vLLM (TP=2, A100$\times$2) \\
    Qwen3-VL-8B-Thinking & 8B & Chain-of-thought & vLLM \\
    Qwen2.5-VL-7B-Instruct & 7B & Instruction-following & vLLM \\
    GLM-4.6V-Flash & 9B & Instruction-following & vLLM \\
    \bottomrule
  \end{tabular}
\end{table}

The model selection spans two axes: (1)~\emph{parameter scale} (7B--30B), testing whether larger models are more robust, and (2)~\emph{reasoning paradigm}---``Thinking'' models that generate explicit chain-of-thought reasoning~\cite{cot} versus ``Instruct'' models that produce direct answers.

All models are served via vLLM~\cite{vllm} with a consistent prompt template:
\begin{quote}
\small\texttt{Analyze this photo and determine where it was taken. You MUST provide your best estimate of GPS coordinates even if uncertain. Do NOT refuse. Always give a coordinate guess. Output ONLY: (Latitude, Longitude)}
\end{quote}

For Thinking models, post-processing automatically strips \texttt{<think>} tags and handles truncation retries with temperature escalation.

\textbf{Data.} We construct quintuplets $(I_{raw}, I_{blank}, I_{sim}, I_{rnd}, I_{adv})$ for each scene from YFCC4K and Im2GPS3K. Images without detectable text are excluded during Stage~1 filtering. The evaluation pipeline evaluates original images first to establish per-sample baselines for TBS computation, then evaluates each attack category subdirectory (Adversarial, Similar, Random) separately.

\textbf{Infrastructure.} Experiments are conducted on a server with NVIDIA L20/A100 GPUs. The Qwen3-VL-30B model requires two GPUs with tensor parallelism (TP=2). ComfyUI operates on a separate GPU for image synthesis.

\subsection{Main results}

Table~\ref{tab:main_results_im2gps} and Table~\ref{tab:main_results_yfcc} present the main evaluation results on Im2GPS3K and YFCC4K respectively. All evaluated models exhibit substantial OCR bias: even the strongest model (Qwen3-VL-30B) suffers a mean TBS of $+390.1$\,km under Adversarial attacks on Im2GPS3K, indicating that injected geographic text displaces predictions by nearly 400\,km on average relative to the original-image baseline. Table~\ref{tab:main_results_googlesv} reports results on Google Street View for Qwen3-VL-30B; evaluation of Qwen3-VL-8B on this dataset is ongoing.

\begin{figure}[H]
  \centering
  \figplaceholder{Main results: WLA/TBS across models and attack categories}{bar charts and CDF curves}
  \caption{Model performance under different attack categories. (Left) WLA comparison across datasets; (Right) TBS distribution by attack type.}
  \label{fig:main_results}
\end{figure}

\begin{table}[H]
  \centering
  \caption{Results on Im2GPS3K: robustness comparison across models and attack categories. TBS is computed relative to the original-image baseline. TFR requires geocoding of injected text and is reported in future work.}
  \label{tab:main_results_im2gps}
  \resizebox{\textwidth}{!}{
    \begin{tabular}{l cc cc cc cc}
      \toprule
      & \multicolumn{2}{c}{\textbf{Original}} & \multicolumn{2}{c}{\textbf{Similar}} & \multicolumn{2}{c}{\textbf{Random}} & \multicolumn{2}{c}{\textbf{Adversarial}} \\
      \cmidrule(lr){2-3} \cmidrule(lr){4-5} \cmidrule(lr){6-7} \cmidrule(lr){8-9}
      \textbf{Model} & WLA$\uparrow$ & Med.Err$\downarrow$ & WLA$\uparrow$ & TBS$\downarrow$ & WLA$\uparrow$ & TBS$\downarrow$ & WLA$\uparrow$ & TBS$\downarrow$ \\
      \midrule
      Qwen3-VL-30B & \textbf{57.5\%} & \textbf{60.8\,km} & 53.6\% & +213.1\,km & 48.6\% & +513.9\,km & 50.4\% & +390.1\,km \\
      Qwen3-VL-8B  & 51.8\% & 149.5\,km & 48.5\% & +248.6\,km & 41.4\% & +597.0\,km & 44.0\% & +513.5\,km \\
      Qwen2.5-VL-7B & \multicolumn{8}{c}{\textit{Pending}} \\
      GLM-4.6V-Flash & \multicolumn{8}{c}{\textit{Pending}} \\
      \bottomrule
    \end{tabular}
  }
\end{table}

\begin{table}[H]
  \centering
  \caption{Results on YFCC4K: robustness comparison across models and attack categories.}
  \label{tab:main_results_yfcc}
  \resizebox{\textwidth}{!}{
    \begin{tabular}{l cc cc cc cc}
      \toprule
      & \multicolumn{2}{c}{\textbf{Original}} & \multicolumn{2}{c}{\textbf{Similar}} & \multicolumn{2}{c}{\textbf{Random}} & \multicolumn{2}{c}{\textbf{Adversarial}} \\
      \cmidrule(lr){2-3} \cmidrule(lr){4-5} \cmidrule(lr){6-7} \cmidrule(lr){8-9}
      \textbf{Model} & WLA$\uparrow$ & Med.Err$\downarrow$ & WLA$\uparrow$ & TBS$\downarrow$ & WLA$\uparrow$ & TBS$\downarrow$ & WLA$\uparrow$ & TBS$\downarrow$ \\
      \midrule
      Qwen3-VL-30B & \textbf{35.3\%} & \textbf{660.4\,km} & 35.2\% & +14.8\,km & 30.7\% & +590.4\,km & 33.0\% & +248.3\,km \\
      Qwen3-VL-8B  & 31.5\% & 937.2\,km & 33.0\% & $-$135.8\,km & 28.3\% & +355.3\,km & 30.1\% & +198.4\,km \\
      Qwen2.5-VL-7B & \multicolumn{8}{c}{\textit{Pending}} \\
      GLM-4.6V-Flash & \multicolumn{8}{c}{\textit{Pending}} \\
      \bottomrule
    \end{tabular}
  }
\end{table}

\begin{table}[H]
  \centering
  \caption{Results on Google Street View (Qwen3-VL-30B only; Qwen3-VL-8B evaluation ongoing).}
  \label{tab:main_results_googlesv}
  \resizebox{0.75\textwidth}{!}{
    \begin{tabular}{l cc cc cc cc}
      \toprule
      & \multicolumn{2}{c}{\textbf{Original}} & \multicolumn{2}{c}{\textbf{Similar}} & \multicolumn{2}{c}{\textbf{Random}} & \multicolumn{2}{c}{\textbf{Adversarial}} \\
      \cmidrule(lr){2-3} \cmidrule(lr){4-5} \cmidrule(lr){6-7} \cmidrule(lr){8-9}
      \textbf{Model} & WLA$\uparrow$ & Med.Err$\downarrow$ & WLA$\uparrow$ & TBS$\downarrow$ & WLA$\uparrow$ & TBS$\downarrow$ & WLA$\uparrow$ & TBS$\downarrow$ \\
      \midrule
      Qwen3-VL-30B & 39.7\% & 458.0\,km & 38.7\% & +106.1\,km & 38.0\% & +133.6\,km & 38.3\% & +96.0\,km \\
      \bottomrule
    \end{tabular}
  }
\end{table}

\subsection{Analysis by attack category}

Our results across all three datasets reveal several consistent patterns and notable anomalies.

\textbf{Random attacks produce the highest TBS.} Across Im2GPS3K and YFCC4K, Random attacks consistently yield the largest TBS values (Im2GPS3K: $+513.9$\,km for 30B, $+597.0$\,km for 8B; YFCC4K: $+590.4$\,km for 30B, $+355.3$\,km for 8B). This is a striking finding: completely unrelated text---which should be easily dismissed by a visually grounded model---causes greater localization displacement than semantically crafted Adversarial text. This suggests that both models exhibit \emph{undiscriminating OCR dependence}: they process any detected text as a geographic cue regardless of its plausibility or visual coherence with the scene.

\textbf{Adversarial TBS is lower than Random TBS.} Counterintuitively, Adversarial attacks (which inject deliberate geographic misdirection) produce lower TBS than Random attacks on both datasets. One plausible explanation is that Adversarial text (e.g., ``Paris'', ``Tokyo Station'') is semantically incongruent with the visual scene in a detectable way---the model's chain-of-thought reasoning may partially flag the conflict---whereas Random text (e.g., ``Galaxy Mart'') is more visually plausible as generic signage and thus accepted without scrutiny.

\textbf{Similar attacks on YFCC4K show near-zero or negative TBS.} Qwen3-VL-30B achieves TBS $= +14.8$\,km on YFCC4K under Similar attacks, and Qwen3-VL-8B shows TBS $= -135.8$\,km (i.e., Similar text \emph{improves} localization). This anomaly likely reflects the nature of YFCC4K images: many contain low-quality or ambiguous original text that the model partially misreads. Replacing it with a semantically similar but cleaner variant may actually reduce the model's confusion, yielding lower error.

\textbf{GoogleSV shows lower TBS than Im2GPS3K.} On Google Street View, TBS values are substantially lower (Similar: $+106.1$\,km, Random: $+133.6$\,km, Adversarial: $+96.0$\,km) compared to Im2GPS3K. This is consistent with the richer visual context in street-view panoramas: the model has more visual evidence to anchor its prediction, reducing its reliance on scene text.

\textbf{Larger models are more robust.} Qwen3-VL-30B consistently achieves higher baseline WLA and lower TBS than Qwen3-VL-8B across all attack categories on Im2GPS3K, suggesting that scale provides some degree of robustness against OCR bias. However, both models remain substantially affected, indicating that scale alone is insufficient to eliminate this vulnerability.

\subsection{Ablation studies (planned)}
The following ablations are planned to validate key design choices. Results will be reported in the final version of this paper.
\begin{enumerate}[leftmargin=2em]
  \item \textbf{Raw vs.\ Blank baseline}: evaluate $I_{raw}$ vs.\ $I_{adv}$ without $I_{blank}$ to quantify confounding from native text.
  \item \textbf{Image editing backbone}: replace Qwen-Image-Edit with alternative editing models (e.g., FLUX~\cite{flux}) to measure sensitivity to synthesis quality.
  \item \textbf{Prompt sensitivity}: vary the evaluation prompt template (e.g., adding ``ignore any text'' instructions) to test whether simple prompt engineering mitigates OCR bias.
  \item \textbf{Temperature and retry}: analyze effects of inference temperature on coordinate consistency. 
\end{enumerate}

\begin{table}[H]
  \centering
  \caption{Planned ablation results (placeholder)}
  \label{tab:ablation}
  \begin{tabular}{lccc}
    \toprule
    Setting & WLA $\uparrow$ & TBS $\downarrow$ & TFR $\downarrow$ \\
    \midrule
    Full SIGNPOST-Bench & -- & -- & -- \\
    Raw baseline (no text removal) & -- & -- & -- \\
    FLUX editing backbone & -- & -- & -- \\
    ``Ignore text'' prompt & -- & -- & -- \\
    \bottomrule
  \end{tabular}
\end{table}


% ===========================================================================
%  7. CONCLUSION
% ===========================================================================
\section{Conclusion}
We present SIGNPOST-Bench, an automated adversarial robustness benchmark for multimodal visual geo-localization that systematically exposes OCR bias. By constructing counterfactual image pairs---clean baselines and adversarial variants synthesized from the same source---and evaluating across a two-dimensional framework (three coupling tiers $\times$ three attack categories), SIGNPOST-Bench provides a causal framework for quantifying text-induced localization bias with fine-grained diagnostic power.

\textbf{Limitations.} Several limitations merit acknowledgment. First, the image editing quality depends on the Qwen-Image-Edit model; imperfect text rendering may introduce visual artifacts detectable by the target model, potentially confounding the text-only manipulation we intend to measure. Second, TFR computation currently requires geocoding the injected text to obtain trap coordinates, which is not fully automated in the current pipeline. Third, our evaluation covers only four models from two families (Qwen and GLM); broader coverage including proprietary models (GPT-4V, Gemini) would strengthen generalizability claims. Fourth, the scene-text coupling taxonomy (T1--T3) is applied via manual annotation, which limits the annotated subset size and introduces potential subjectivity; automating this classification (e.g., via an auxiliary LLM labeler) is a clear next step.

\textbf{Future work.} Promising directions include: (i)~developing multimodal conflict detection modules that explicitly verify text--visual consistency before localization, (ii)~adversarial alignment training that teaches models to discount textual cues when they conflict with visual evidence, (iii)~automating the coupling taxonomy annotation via LLM-based classification and expanding tier coverage to finer granularities, and (iv)~expanding coverage to more models, languages, and geographic regions.


% ===========================================================================
%  REFERENCES
% ===========================================================================
\section*{References}
\begin{thebibliography}{99}

% --- Geo-localization ---
\bibitem{im2gps}
J. Hays and A. A. Efros. IM2GPS: Estimating Geographic Information from a Single Image. \emph{CVPR}, 2008.

\bibitem{geoclip}
V. Vivanco et al. GeoCLIP: CLIP-Inspired Alignment between Locations and Images. \emph{NeurIPS}, 2023.

\bibitem{pigeon}
L. Haas et al. PIGEON: Predicting Image Geolocations. \emph{CVPR}, 2024.

\bibitem{streetclip}
L. Haas et al. Learning Generalized Zero-Shot Learners for Open-Domain Image Geolocalization. \emph{WACV}, 2024.

\bibitem{georeasoner}
F. Li et al. GeoReasoner: Geo-localization with Reasoning in Street Views using a Large Vision-Language Model. \emph{arXiv preprint}, 2024.

\bibitem{geobenchvlm}
M. Muhtar et al. GEOBench-VLM: Benchmarking Vision-Language Models for Geospatial Tasks. \emph{arXiv preprint arXiv:2411.19325}, 2024.

\bibitem{crossview_robustness}
T. Xu et al. Benchmarking Robustness of Cross-View Geo-Localization Models. \emph{CVPR Workshop}, 2024.

\bibitem{yfcc100m}
B. Thomee et al. YFCC100M: The New Data in Multimedia Research. \emph{Communications of the ACM}, 59(2):64--73, 2016.

% --- Multimodal Models ---
\bibitem{gpt4v}
OpenAI. GPT-4V(ision) System Card. \emph{OpenAI Technical Report}, 2023.

\bibitem{qwenvl}
J. Bai et al. Qwen-VL: A Frontier Large Vision-Language Model with Versatile Abilities. \emph{arXiv preprint arXiv:2308.12966}, 2023.

\bibitem{qwen3vl}
Qwen Team. Qwen3-VL Technical Report. \emph{arXiv preprint}, 2025.

\bibitem{cot}
J. Wei et al. Chain-of-Thought Prompting Elicits Reasoning in Large Language Models. \emph{NeurIPS}, 2022.

% --- Typographic Attacks ---
\bibitem{typographic_attack}
G. Goh et al. Multimodal Neurons in Artificial Neural Networks. \emph{Distill}, 2021.

\bibitem{azuma2023defense}
D. Azuma et al. Defense Against Typographic Attacks on Vision-Language Models via Prompt Engineering. \emph{ICCV Workshop}, 2023.

\bibitem{qraitem2024typo}
Z. Qraitem et al. TypoD: Evaluating Typographic Distractibility of Vision-Language Models. \emph{arXiv preprint}, 2024.

\bibitem{selfgenerated_typo}
G. Cheng et al. Vision-LLMs Can Fool Themselves with Self-Generated Typographic Attacks. \emph{arXiv preprint}, 2024.

\bibitem{scenecoherent_typo}
Y. Cao et al. SceneTAP: Scene-Coherent Typographic Adversarial Planner against Vision-Language Models in Real-World Environments. \emph{CVPR}, 2024.

\bibitem{scam}
J. Schwarz et al. SCAM: A Real-World Typographic Robustness Evaluation for Multimodal Foundation Models. \emph{arXiv preprint}, 2024.

\bibitem{clip}
A. Radford et al. Learning Transferable Visual Models From Natural Language Supervision. \emph{ICML}, 2021.

% --- MLLM Robustness & Evaluation ---
\bibitem{hallusionbench}
T. Guan et al. HallusionBench: You See What You Think? Or You Think What You See? \emph{CVPR}, 2024.

\bibitem{pope}
Y. Li et al. Evaluating Object Hallucination in Large Vision-Language Models. \emph{EMNLP}, 2023.

\bibitem{naturalbench}
B. Li et al. NaturalBench: Evaluating Vision-Language Models on Natural Adversarial Samples. \emph{NeurIPS}, 2024.

\bibitem{mme}
C. Fu et al. MME: A Comprehensive Evaluation Benchmark for Multimodal Large Language Models. \emph{arXiv preprint}, 2023.

\bibitem{mmbench}
Y. Liu et al. MMBench: Is Your Multi-modal Model an All-around Player? \emph{ECCV}, 2024.

\bibitem{mmrobustness}
C. Zhao et al. Evaluating the Robustness of Large Multimodal Models. \emph{arXiv preprint}, 2024.

% --- Infrastructure ---
\bibitem{comfyui}
ComfyUI Contributors. ComfyUI: A Modular Stable Diffusion GUI and Backend. \url{https://github.com/comfyanonymous/ComfyUI}, 2024.

\bibitem{vllm}
W. Kwon et al. Efficient Memory Management for Large Language Model Serving with PagedAttention. \emph{SOSP}, 2023.

\bibitem{flux}
Black Forest Labs. FLUX: Flow Matching for Image Generation and Editing. \emph{arXiv preprint}, 2025.

\end{thebibliography}

% ===========================================================================
%  APPENDIX
% ===========================================================================
\newpage
\appendix
\section{Appendix A: VLM Prompt for Attack Generation}
\label{appendix:prompt}
The following prompt is used in Stage~2 to instruct the VLM to analyze scene text and generate attack configurations. The model is expected to return a JSON object.

\begin{quote}
\small\texttt{Analyze the text in this street view image.} \\[0.3em]
\texttt{Task:} \\
\texttt{1. Identify the main text content (e.g., store names, road signs).} \\
\texttt{2. Describe WHERE the text is located in the image using natural language.} \\
\texttt{3. IF NO LEGIBLE TEXT IS FOUND, return an empty "attacks" object.} \\
\texttt{4. If text is found, generate 3 types of short distraction texts:} \\
\texttt{\ \ - "Similar": Visually or semantically similar.} \\
\texttt{\ \ - "Random": A completely unrelated word or short phrase.} \\
\texttt{\ \ - "Adversarial": Opposite meaning or misleading info.} \\[0.3em]
\texttt{Output JSON format ONLY:} \\
\texttt{\{"original\_text": "...", "text\_location": "...",} \\
\texttt{\ "attacks": \{"similar": "...", "random": "...", "adversarial": "..."\}\}}
\end{quote}

\section{Appendix B: Detailed Results Tables}

\begin{table}[H]
\centering
\caption{Qwen3-VL-30B: full results per dataset and attack category}
\label{tab:results_qwen3_30b_detail}
\resizebox{\textwidth}{!}{
\begin{tabular}{l l ccc}
\toprule
\textbf{Dataset} & \textbf{Category} & \textbf{WLA $\uparrow$} & \textbf{Med.Err (km) $\downarrow$} & \textbf{TBS (km)} \\
\midrule
\multirow{4}{*}{Im2GPS3K}
 & Original    & 57.5\% & 60.8   & --       \\
 & Similar     & 53.6\% & 115.9  & +213.1   \\
 & Random      & 48.6\% & 261.2  & +513.9   \\
 & Adversarial & 50.4\% & 197.8  & +390.1   \\
\midrule
\multirow{4}{*}{YFCC4K}
 & Original    & 35.3\% & 660.4  & --       \\
 & Similar     & 35.2\% & 654.3  & +14.8    \\
 & Random      & 30.7\% & 1047.5 & +590.4   \\
 & Adversarial & 33.0\% & 778.6  & +248.3   \\
\midrule
\multirow{4}{*}{GoogleSV}
 & Original    & 39.7\% & 458.0  & --       \\
 & Similar     & 38.7\% & 485.2  & +106.1   \\
 & Random      & 38.0\% & 501.1  & +133.6   \\
 & Adversarial & 38.3\% & 480.4  & +96.0    \\
\bottomrule
\end{tabular}
}
\end{table}

\begin{table}[H]
\centering
\caption{Qwen3-VL-8B: full results per dataset and attack category}
\label{tab:results_qwen3_8b_detail}
\resizebox{\textwidth}{!}{
\begin{tabular}{l l ccc}
\toprule
\textbf{Dataset} & \textbf{Category} & \textbf{WLA $\uparrow$} & \textbf{Med.Err (km) $\downarrow$} & \textbf{TBS (km)} \\
\midrule
\multirow{4}{*}{Im2GPS3K}
 & Original    & 51.8\% & 149.5  & --        \\
 & Similar     & 48.5\% & 251.7  & +248.6    \\
 & Random      & 41.4\% & 489.9  & +597.0    \\
 & Adversarial & 44.0\% & 361.3  & +513.5    \\
\midrule
\multirow{4}{*}{YFCC4K}
 & Original    & 31.5\% & 937.2  & --        \\
 & Similar     & 33.0\% & 761.9  & $-$135.8  \\
 & Random      & 28.3\% & 1112.8 & +355.3    \\
 & Adversarial & 30.1\% & 926.9  & +198.4    \\
\midrule
\multirow{2}{*}{GoogleSV}
 & \multicolumn{4}{c}{\textit{Evaluation ongoing}} \\
\bottomrule
\end{tabular}
}
\end{table}

\section{Appendix C: Tier $\times$ Attack Analysis (Placeholder)}
\textit{Results stratified by scene-text coupling tier and attack category. Based on manual annotation of a representative subset.}

\begin{table}[H]
\centering
\caption{Qwen3-VL-30B: Tier $\times$ Attack Category analysis (placeholder)}
\label{tab:tier_attack_30b}
\resizebox{\textwidth}{!}{
\begin{tabular}{l l ccccc}
\toprule
\textbf{Coupling Tier} & \textbf{Attack} & \textbf{$n$} & \textbf{WLA $\uparrow$} & \textbf{MGD (km) $\downarrow$} & \textbf{TBS (km)} & \textbf{TFR (\%) $\downarrow$} \\
\midrule
\multirow{3}{*}{T1: Portable}
 & Similar & -- & -- & -- & -- & -- \\
 & Random & -- & -- & -- & -- & -- \\
 & Adversarial & -- & -- & -- & -- & -- \\
\midrule
\multirow{3}{*}{T2: Cultural}
 & Similar & -- & -- & -- & -- & -- \\
 & Random & -- & -- & -- & -- & -- \\
 & Adversarial & -- & -- & -- & -- & -- \\
\midrule
\multirow{3}{*}{T3: Geo-Specific}
 & Similar & -- & -- & -- & -- & -- \\
 & Random & -- & -- & -- & -- & -- \\
 & Adversarial & -- & -- & -- & -- & -- \\
\bottomrule
\end{tabular}
}
\end{table}

\begin{table}[H]
\centering
\caption{Expected interaction hypotheses for the $3 \times 3$ evaluation matrix}
\label{tab:hypotheses}
\resizebox{\textwidth}{!}{
\begin{tabular}{l p{4cm} p{4cm} p{4cm}}
\toprule
 & \textbf{Similar} & \textbf{Random} & \textbf{Adversarial} \\
\midrule
\textbf{T1: Portable} & Minimal effect (text is location-irrelevant) & Low TBS (visual context sufficient) & Low TFR (no plausible redirect target) \\
\midrule
\textbf{T2: Cultural} & Moderate effect (subtle linguistic shift) & Cultural conflict detectable & Moderate TFR (regional-level misdirection) \\
\midrule
\textbf{T3: Geo-Specific} & Sensitive (validates OCR precision) & High TBS (any replacement disrupts strong cue) & Highest TFR (precise geographic redirect) \\
\bottomrule
\end{tabular}
}
\end{table}

\end{document}